\begin{theorem}[Interior support of coherent regular torus closures]%
    \label{thm:peps_regular_torus_sector_interior_support}
    \lean{TNLean.PEPS.range_torusClosureSuperpositionCut_le_regionGroundSpace_interior_rectangle,
        TNLean.PEPS.IsGInjective.range_torusClosureSuperpositionCut_eq_regionGroundSpace_interior_rectangle,
        TNLean.PEPS.IsGInjective.range_torusClosureReducedDensity_eq_regionGroundSpace_interior_rectangle}
    \leanok
    \uses{thm:peps_torus_closure_interior_parent_constraints,
        thm:peps_regular_torus_rectangle_entropy_bound}
    Let a native four-leg regular $G$-injective tensor be placed on a
    square torus with both periods at least three. Let
    $\Psi=\sum_i\mu_i\Psi_{g_i,h_i}$ be a nonzero coherent sum of
    closure vectors, where $g_i h_i=h_i g_i$ for every $i$.
    Let $R$ be a positive coordinate rectangle whose starting
    coordinates are positive and whose ending coordinates are
    strictly below the torus periods. Then the physical coefficient
    matrix across $R$ has range exactly the original regional
    ground space $\mathcal G_R$. Equivalently, the actual reduced
    density of $\Psi$ satisfies
    \[
        \operatorname{ran}\rho_R=\mathcal G_R.
    \]
    This is an interior support consequence of
    \cite[Theorem~5.5 and Corollary~6.10]{Schuch2010PEPS}; it does
    not classify the full torus parent-Hamiltonian ground space.
\end{theorem}

\begin{proof}\leanok
    \uses{thm:peps_torus_closure_interior_parent_constraints,
        lem:peps_torus_rectangle_boundary_card,
        thm:peps_regular_region_dimension}
    Every complementary physical slice of $\Psi$ lies in
    $\mathcal G_R$, so the coefficient range is contained in that
    space. The simply connected regular Schmidt-rank theorem gives
    rank $|G|^{|\partial R|-1}$ for the nonzero coherent sum.
    The connected open-region contraction has precisely the same
    dimension. Inclusion and equality of dimensions therefore give
    equality of the two spaces. The reduced density is the product
    of the coefficient matrix and its adjoint, which has the same
    range as the coefficient matrix.
\end{proof}
